\section{Independent Context-Transfer Trial}

\subsection{Research objective}
The prospective experiment tests whether practice of a common executive-control rule across several interaction contexts improves an independent transfer outcome relative to repeated practice in one context. This is a proposed protocol, not a report of observed effects.

\subsection{Randomized comparison}
Children meeting prespecified eligibility and assent criteria will be assigned to fixed-context or varied-context training. Both arms will share cognitive rules, challenge constraints, reward logic, planned training dose, and stopping procedures. The varied-context arm will use stationary, embodied non-AR, and embodied AR tasks on technically qualified family-owned Android devices. Hardware eligibility will be determined before randomization. Device failure, refusal, and missed sessions will be recorded rather than silently excluded.

\subsection{Independent endpoint}
The primary assessment must use novel stimuli and an untrained response format outside the AR interface. Its age suitability, scoring, reliability, timing, and minimally useful difference will be established before recruitment. Baseline, post-training, and retention assessments are planned; the primary time point must be selected in advance. Training accuracy, movement exposure, enjoyment, and recovery telemetry are secondary and cannot substitute for independent transfer.

\subsection{Primary estimand}
Let \(Y_i(z)\) be the primary post-training transfer score that child \(i\) would have under assignment \(z\), with \(z=1\) for varied-context training and \(z=0\) for fixed-context training. The assignment-based target is
\[
\Delta_T=\mathbb{E}[Y_i(1)-Y_i(0)].
\]
A baseline-adjusted group comparison will estimate this effect. Prespecified covariates, uncertainty intervals, and the decision threshold will be frozen before outcome access. Retention and change-score analyses are secondary. Postrandomization adherence will not be used as a baseline adjustment.

\subsection{Missingness and interpretation}
Withdrawal, safety stopping, hardware failure, and missing transfer assessments will be enumerated by randomized arm. Training discontinuation does not automatically remove a child from the assignment-based analysis. Missing outcomes require sensitivity analyses allowing departures from missing-at-random assumptions. A gain restricted to practiced tasks does not establish transfer. Any transfer finding remains distinct from clinical ADHD symptom improvement, which requires separate validated outcomes.